% Generated draft: requires mathematical and source review.
\documentclass[11pt]{article}
\usepackage[margin=25mm]{geometry}
\usepackage{fontspec}
\setmainfont{Latin Modern Roman}
\usepackage{amsmath,amssymb,mathtools,unicode-math}
\setmathfont{Latin Modern Math}
\usepackage{enumitem,needspace,xurl,hyperref}
\setlength{\parindent}{0pt}
\setlength{\parskip}{5pt}
\setlength{\emergencystretch}{3em}
\newcommand{\N}{\mathbb{N}}
\newcommand{\Z}{\mathbb{Z}}
\newcommand{\Q}{\mathbb{Q}}
\newcommand{\R}{\mathbb{R}}
\newcommand{\C}{\mathbb{C}}
\newcommand{\F}{\mathbb{F}}
\newcommand{\A}{\mathbb{A}}
\newcommand{\PP}{\mathbb{P}}
\newcommand{\E}{\mathbb{E}}
\newcommand{\eps}{\varepsilon}
\newcommand{\dd}{\,\mathrm{d}}
\begin{document}
\section*{UnsolvedMath ID: 30006523}
\textbf{Source problem code:} catalogue entry
\subsection{Definitions and mathematical statement}
Use the conventional ordinary single Grothendieck polynomials and ordinary Schubert polynomials, with the normalization specified below. The supplied excerpt does not explicitly define these polynomial conventions or its Schubert-support notation; this interpretation requires confirmation against the report.

Let \(n\geq 1\), let \(S_n\) be the symmetric group on \(\{1,\ldots,n\}\), and let \(s_i=(i,i+1)\) for \(1\leq i<n\). Write \(\ell(w)\) for the number of inversions of \(w\), and let \(w_0(i)=n+1-i\). For a polynomial \(f\in\Z[x_1,\ldots,x_n]\), let \(s_i f\) interchange \(x_i\) and \(x_{i+1}\), and define
\[
\partial_i f=\frac{f-s_i f}{x_i-x_{i+1}},
\qquad
\pi_i f=\partial_i\bigl((1-x_{i+1})f\bigr).
\]
The Schubert polynomials \(\mathfrak S_w\) and Grothendieck polynomials \(G_w\) are normalized by
\[
\mathfrak S_{w_0}=G_{w_0}=\prod_{j=1}^{n-1}x_j^{n-j},
\]
and, whenever \(\ell(ws_i)=\ell(w)+1\), by
\[
\mathfrak S_w=\partial_i\mathfrak S_{ws_i},
\qquad
G_w=\pi_iG_{ws_i}.
\]
Write the Schubert expansion as
\[
G_w=\sum_{u\in S_n}c_{w,u}\mathfrak S_u,
\qquad c_{w,u}\in\Z,
\]
and define its Schubert support by
\[
\operatorname{Supp}_S(G_w)=\{u\in S_n:c_{w,u}\neq0\}.
\]
Here support records nonzero coefficients, regardless of sign.

The left weak order is defined by
\[
u\leq_L v
\quad\Longleftrightarrow\quad
\text{there exists }a\in S_n\text{ such that }v=au
\text{ and }\ell(v)=\ell(a)+\ell(u).
\]
For \(u,u'\in S_n\), put
\[
[u,u']_L=\{v\in S_n:u\leq_L v\leq_L u'\}.
\]

\textbf{Conjecture 3 in the supplied source.} For every integer \(n\geq1\) and every \(w\in S_n\), the set \(\operatorname{Supp}_S(G_w)\) is order-convex in left weak order. Explicitly,
\[
\forall u,u',v\in S_n,\qquad
\bigl(c_{w,u}\neq0\ \land\ c_{w,u'}\neq0\ \land\ u\leq_L v\leq_L u'\bigr)
\ \Longrightarrow\ c_{w,v}\neq0.
\]
Equivalently, every interval with both endpoints in the support is contained in the support. For incomparable endpoints the interval is empty, so the condition imposes no additional requirement.

\textbf{Supplied qualifications.} The source excerpt reports that these conjectures are resolved for Grassmannian permutations and for permutations of the form \(1\times a\), where \(a\) is dominant. Under the standard terminology, Grassmannian means having at most one descent, dominant means avoiding the pattern \(132\), and, for \(a\in S_m\), the permutation \(1\times a\in S_{m+1}\) satisfies \((1\times a)(1)=1\) and \((1\times a)(j+1)=a(j)+1\). The explicit dominant-related case in the excerpt is thus \(1\times a\); it should not silently be enlarged to an unrestricted claim about dominant permutations.

The excerpt separately reports results for inverse fireworks permutations only for its first two conjectures. It does not assert that those results establish the present Conjecture 3. It also states that Conjectures 1 and 2 would imply analogous conjectures in its reference \([3]\), replacing Schubert support by monomial support and left weak order by divisibility. Those two conjectures and the identity of \([3]\) are not supplied here, so this contextual implication cannot be formulated more precisely from the record.
\subsection{Short English statement}
For every permutation \(w\), expand its Grothendieck polynomial in Schubert polynomials and retain the permutations with nonzero coefficients. The question asks whether every permutation lying between two retained permutations in left weak order must also be retained. The supplied record labels the problem partially solved and reports particular established cases; these reports and the polynomial conventions require human review.
\subsection{Sources}
\begin{itemize}
\item Local UnsolvedMath registry: supplied complete record with permanent ID \(30006523\), title \emph{Weak-Order Convexity of Grothendieck Schubert Supports}, and exact source code \texttt{catalogue entry}. The record assigns proposed year \(2026\), but names no individual proposer.
\item Supplied primary report citation: \emph{Enumerative Combinatorics} (2026), Oberwolfach Reports, DOI \texttt{10.4171/owr/2026/2}, pp.\ 81--156. \url{https://doi.org/10.4171/owr/2026/2}. The supplied excerpt labels the target Conjecture 3; no more specific page or contribution author is given.
\item Mészáros--Setiabrata--St. Dizier, \emph{On the support of Grothendieck polynomials}. Supplied literature link: \url{https://arxiv.org/abs/2201.09452}.
\item Chou--Yu, \emph{Grothendieck polynomials of inverse fireworks permutations}. Supplied literature link: \url{https://arxiv.org/abs/2403.12186}.
\item Dou--Fan--Liu, \emph{Lattice-free Schubitopes}. Supplied literature link: \url{https://arxiv.org/abs/2605.10016}.
\item Supplied literature assessment, checked \texttt{2026-08-22}: status \texttt{partially\_solved}. The assessment reports convexity in the listed Grassmannian and dominant cases, with further partial results, but not for all permutations. This is supplied dated curation derived from the OWR record, not an independently verified literature assessment. The record itself requests independent verification of time-sensitive status and cited claims.
\item The supplied registry labels the statement \texttt{corrected\_verified} and explains that only formula delimiters, notation, grammar, and typography were repaired while preserving mathematical content and scope. This is a supplied editorial label, not independent mathematical verification of this draft.
\end{itemize}
\end{document}
